\newpage
\section{CPU Backend}

The CPU backend serves as the reference implementation and fallback for all other backends. It is written entirely in Python with NumPy, ensuring broad compatibility and zero external dependencies beyond the core scientific stack. Despite its simplicity, it incorporates several optimizations that make it suitable for moderate-scale continuous field operations.

\subsection{Vectorized Field Operations}

All core operations are implemented using NumPy's vectorized primitives, avoiding explicit Python loops. The \texttt{create\_continuous\_field} method generates a random field and normalizes it:

\begin{lstlisting}[language=Python]
def create_continuous_field(self, dimensions: int) -> np.ndarray:
    field = np.random.randn(dimensions).astype(self.precision)
    field = self._normalize(field)
    # ... caching logic ...
    return field
\end{lstlisting}

The \texttt{field\_update} method performs a single time step using element-wise array operations. It simulates analog continuity by adding small Gaussian noise scaled by \(\sqrt{\Delta t}\), a discretization of the continuous-time Langevin equation:

\begin{equation}
d\mathbf{x}_t = -\eta \nabla \mathcal{L}(\mathbf{x}_t)dt + \sigma d\mathbf{W}_t \quad\Rightarrow\quad \mathbf{x}_{t+1} \approx \mathbf{x}_t + \sigma \sqrt{\Delta t}\,\boldsymbol{\xi}_t
\end{equation}

where \(\boldsymbol{\xi}_t \sim \mathcal{N}(0,I)\) and we omit the drift term for simplicity. The noise scale \(\sigma = 0.01\) is chosen to balance exploration and stability.

\subsection{O(1) Field Identification via Hashing}

A recurring pattern in the APEIRON framework is the need to track the identity of a field across multiple operations. Naively, this would require \(O(N)\) linear scans to locate a field's internal ID. The CPU backend eliminates this overhead by maintaining a \texttt{field\_hash} dictionary that maps an MD5 digest of the field's byte representation to its internal ID. The \texttt{get\_field\_id} method first attempts an \(O(1)\) hash lookup; only on a cache miss does it fall back to a linear scan (and subsequently updates the hash):

\begin{lstlisting}[language=Python]
def get_field_id(self, field: np.ndarray) -> Optional[int]:
    try:
        field_hash = hashlib.md5(field.tobytes()).hexdigest()
        if field_hash in self.field_hash:
            self.metrics['cache_hits'] += 1
            return self.field_hash[field_hash]
    except Exception:
        pass
    self.metrics['cache_misses'] += 1
    # linear scan and hash update
\end{lstlisting}

This optimization is critical for operations like \texttt{find\_stable\_patterns}, which may perform \(O(N^2)\) pairwise comparisons and would otherwise be dominated by ID lookups.

\subsection{Metrics and Observability}

The CPU backend maintains an extensive set of performance metrics, including total fields created, number of updates, interference calculations, pattern searches, and coherence measurements. All timing metrics are updated using an exponential moving average (EMA) to provide stable, real-time estimates of operation latency. These metrics are exposed via the \texttt{get\_metrics()} method and can be scraped by Prometheus when the optional metrics server is enabled.

\subsection{Configurable Precision and Threading}

The backend supports both \texttt{float32} and \texttt{float64} precision via the \texttt{precision} configuration key. Multi-threading is controlled by setting the \texttt{OMP\_NUM\_THREADS} and \texttt{MKL\_NUM\_THREADS} environment variables, allowing the backend to leverage multi-core CPUs without code changes.